% ==============================================================================
% preamble.inc.tex — Преамбула для ВКР (дополнение к классу G7-32)
%
% Класс G7-32 уже загружает: babel, fontspec (XeLaTeX), geometry, fancyhdr,
% indentfirst, caption, amssymb, amsmath, longtable, flafter
% ==============================================================================

\sloppy

% --- Настройки стиля ГОСТ 7-32 ---
% \EqInChapter  % формулы нумеруются в пределах раздела
% \TableInChapter % таблицы нумеруются в пределах раздела
% \PicInChapter   % рисунки нумеруются в пределах раздела

% --- Гипертекстовое оглавление ---
\usepackage[
  bookmarks=true, colorlinks=true, unicode=true,
  urlcolor=black, linkcolor=black, anchorcolor=black,
  citecolor=black, menucolor=black, filecolor=black,
]{hyperref}

\urlstyle{same}

% --- Включение PDF-страниц (для титульного листа) ---
\usepackage{pdfpages}

% --- Графика ---
\usepackage{graphicx}
\graphicspath{{images/},{../latex/images/}}

% --- TikZ ---
\usepackage{tikz}
\usetikzlibrary{arrows,positioning,shadows}

% --- Списки и таблицы ---
\usepackage{enumerate}
\usepackage{multirow}
\usepackage{paralist,array}
\usepackage{tabularx}
\usepackage{dcolumn}
\usepackage{booktabs}

% --- Полуторный интервал ---
\usepackage{setspace}
\setstretch{1.43}

% --- Компактный нумерованный список ---
\newenvironment{compactlist}{%
    \renewcommand{\labelenumi}{\hspace{13mm}\arabic{enumi}.}
    \setlength{\topsep}{0pt}%
    \setlength{\partopsep}{0pt}%
    \setlength{\itemsep}{0pt}%
    \setlength{\parsep}{0pt}%
    \begin{enumerate}%
}{%
    \end{enumerate}%
}

% --- Отступы в оглавлении ---
\makeatletter
\renewcommand*\l@section{\@dottedtocline{1}{0em}{2.3em}}
\renewcommand*\l@subsection{\@dottedtocline{2}{0em}{3.2em}}
\renewcommand*\l@subsubsection{\@dottedtocline{3}{0em}{4.1em}}
\makeatother

% --- Подписи к рисункам и таблицам ---
\usepackage{ragged2e}

\DeclareCaptionJustification{centerednohyphen}{%
  \Centering\hyphenpenalty=10000\exhyphenpenalty=10000%
}
\DeclareCaptionLabelSeparator{shortdash}{ \textendash\ }

% --- Настройки float ---
\usepackage{float}
\setlength{\floatsep}{0pt}
\setlength{\textfloatsep}{10pt}
\setlength{\intextsep}{7pt}

\captionsetup{
  aboveskip=7pt, belowskip=0pt,
  justification=centerednohyphen,
  singlelinecheck=false,
  labelsep=shortdash
}
\captionsetup[table]{justification=RaggedRight, singlelinecheck=false}
\captionsetup[figure]{justification=centerednohyphen, singlelinecheck=false}

% --- Счётчики ---
\usepackage{chngcntr}

% --- Рамки и цвет ---
\usepackage{mdframed}
\usepackage{xcolor}

\mdfsetup{
  linecolor=black,
  linewidth=1pt,
  roundcorner=5pt,
  innerleftmargin=5pt,
  innerrightmargin=5pt,
  innertopmargin=5pt,
  innerbottommargin=5pt,
  skipabove=0pt,
  skipbelow=0pt,
  leftmargin=0pt,
  rightmargin=0pt
}

% --- Листинги кода (minted) ---
% NOTE: Для использования minted необходим Python с pygments.
% Раскомментируйте при необходимости (требует --shell-escape).
% Текущая версия minted (0.6.0) несовместима с Python 3.14+.
%
% \usepackage[chapter]{minted}
%
% \captionsetup[listing]{
%   position=top,
%   justification=RaggedRight,
%   singlelinecheck=false
% }
%
% \floatstyle{plaintop}
% \restylefloat{listing}
%
% \setminted{
%   linenos=true,
%   numberblanklines=false,
%   numbersep=5pt,
%   xleftmargin=1em,
%   fontsize=\small,
%   breakanywhere=true,
%   breakautoindent=true,
%   breaklines=true,
% }
%
% \renewcommand{\listingscaption}{Листинг}
% \renewcommand{\listoflistingscaption}{Листинги}

% --- Листинги кода (listings, альтернативный вариант) ---
\usepackage{listings}
\lstset{
  basicstyle=\fontsize{12pt}{14pt}\selectfont\ttfamily,
  breaklines=true,
  frame=none,
  numbers=left,
  numberstyle=\small,
  tabsize=4,
  showstringspaces=false,
  captionpos=t,
  aboveskip=10pt,
  belowskip=10pt
}
\renewcommand{\lstlistingname}{Листинг}

% --- Оглавление без переносов ---
\makeatletter
\let\oldtableofcontents\tableofcontents
\renewcommand{\tableofcontents}{%
  \begingroup
  \hyphenpenalty=10000
  \exhyphenpenalty=10000
  \oldtableofcontents
  \endgroup
}
\makeatother

% --- Заголовки без висячего отступа ---
\makeatletter
\def\@hangfrom#1{\setbox\@tempboxa\hbox{{#1}}%
      \hangindent 0pt%
      \noindent\box\@tempboxa}
\makeatother

% --- Отступы после окружений ---
\usepackage{etoolbox}
% \AtEndEnvironment{listing}{\vspace{-10pt}} % требует minted
\AtEndEnvironment{table}{\vspace{-10pt}}
\AtEndEnvironment{center}{\vspace{-14pt}}
\AtEndEnvironment{figure}{\vspace{4pt}}

% --- Шрифты (Times New Roman) ---
\setmainfont{Times New Roman}
\newfontfamily\cyrillicfont{Times New Roman}
\setmonofont[Script=Default,Ligatures=TeX]{Courier New}
\newfontfamily\cyrillicfonttt[Script=Cyrillic]{Courier New}

% --- Удалить точку после номера раздела (стиль без точки) ---
\makeatletter
\renewcommand*{\@seccntformat}[1]{%
  \csname the#1\endcsname\hspace{0.5em}%
}
\makeatother

% --- Глубина оглавления (до подразделов) ---
\setcounter{tocdepth}{2}

% --- Пользовательские команды ---
% Сокращение с расшифровкой при первом использовании
\newcommand{\abbr}[2]{#1 (#2)}

% Ссылки на рисунок/таблицу в русском стиле
\newcommand{\figref}[1]{рисунок~\ref{#1}}
\newcommand{\tabref}[1]{таблица~\ref{#1}}
